\chapter{Brain Injuries}
\index{Brain Injuries}

\section*{Definition and Overview}
A brain injury -- more formally a traumatic brain injury (TBI) -- is damage to the brain caused by an external mechanical force, such as a blow to the head, a fall, a road traffic collision, a blast, or a penetrating wound. Brain injuries span an enormous range of severity, from a mild concussion that resolves fully within days or weeks to catastrophic injuries causing prolonged unconsciousness and permanent disability. Two distinct processes are at work in every injury: the primary injury, which is the structural damage done at the moment of impact, and the secondary injury, which is the additional damage that unfolds over the following hours and days from swelling, bleeding, reduced blood flow, and oxygen starvation. Modern management is focused overwhelmingly on preventing and limiting this secondary injury, because it is potentially reversible and is where early, expert care makes the greatest difference to outcome.

\section*{Epidemiology}
Traumatic brain injury is one of the leading causes of death and long-term disability worldwide, and it is the single most important cause of serious neurological injury in people under the age of 45. Global estimates suggest more than 50 million cases per year, with road traffic accidents, falls, and interpersonal violence the dominant mechanisms in most settings. Men are affected roughly twice as often as women. The age distribution is bimodal: falls predominate in young children and in older adults, while motor vehicle collisions and sport-related concussion peak in adolescents and young adults. In older people, even a comparatively minor fall can produce significant injury because the brain is more fragile with age and many older adults take anticoagulant or antiplatelet medication, which markedly increases the risk of bleeding. Mild injuries (concussion) are by far the most common and are frequently unreported or unrecognised.

\section*{Aetiology and Pathophysiology}
Brain injuries arise through several distinct mechanisms:

\begin{itemize}
    \item \textbf{Blunt impact:} the brain is jolted against the inside of the skull, causing bruising (contusion) at the site of impact and often on the opposite side as well (coup--contrecoup injury)
    \item \textbf{Acceleration--deceleration:} rapid deceleration shears and stretches nerve fibres within the brain, producing diffuse axonal injury, which can be severe even without a fracture
    \item \textbf{Penetrating injury:} bullets, shrapnel, or other objects damage brain tissue directly and introduce the risk of infection
    \item \textbf{Skull fracture:} a break in the skull, which may be associated with bleeding around or within the brain
    \item \textbf{Blast injury:} pressure waves from explosions can injure the brain through mechanisms that include the effects of the blast itself
\end{itemize}

The principal types of intracranial bleeding are the epidural haematoma (blood between the skull and the outer brain covering, often after a skull fracture and typically with a lucid interval before deterioration), the subdural haematoma (blood between the brain and its lining, more common in older adults and in those on blood thinners, often accumulating slowly), and intracerebral haemorrhage (bleeding within the brain substance itself). After the primary injury, cerebral oedema, rising intracranial pressure, hypotension, hypoxia, and seizures drive secondary injury, which is why rapid assessment and intensive care are so critical.

\section*{Clinical Features}
Symptoms vary with severity:

\begin{itemize}
    \item \textbf{Mild injury (concussion):} headache, dizziness, nausea, brief confusion or disorientation, short memory gaps, sensitivity to light and noise, difficulty concentrating, and fatigue
    \item \textbf{Moderate to severe injury:} loss of consciousness or a declining level of consciousness, often with behavioural change or agitation
    \item Repeated vomiting, especially after the first few hours
    \item Seizures or fitting
    \item Weakness, numbness, difficulty speaking, or visual disturbance
    \item Unequal pupil size
    \item Blood or clear fluid (cerebrospinal fluid) leaking from the nose or ears
    \item In infants, poor feeding, excessive irritability, or a bulging fontanelle
\end{itemize}

\section*{Differential Diagnosis}
\begin{itemize}
    \item Alcohol or drug intoxication, which can mimic the confusion of brain injury and often coexists with it
    \item Stroke or transient ischaemic attack causing sudden focal symptoms
    \item A seizure disorder -- a fit may cause a fall and a head injury, and a head injury may itself trigger a seizure
    \item Hypoglycaemia, hypoxia, or other metabolic causes of altered consciousness
    \item Meningitis or encephalitis presenting with headache and confusion
    \item Co-existing cervical spine injury, which must be assumed present in any significant head injury until radiologically excluded
\end{itemize}

\section*{When to Seek Medical Care}
Any head injury associated with loss of consciousness, confusion, vomiting, or a severe headache warrants prompt medical assessment. After a mild injury, medical review is advisable if symptoms such as headache, dizziness, or poor concentration are not improving, or if the injured person is elderly, taking blood-thinning medication, has a bleeding tendency, or is a young child or baby.

\textbf{Contact your local emergency services for:}
\begin{itemize}
    \item A decreasing level of consciousness or difficulty waking the person
    \item Repeated vomiting after a head injury
    \item Seizures or fitting
    \item Weakness, numbness, or difficulty speaking
    \item Blood or clear fluid leaking from the nose or ears, or unequal pupils
    \item Any head injury in a baby, in a person taking blood thinners, or with a suspected neck injury
\end{itemize}

\section*{Diagnosis and Investigations}
\begin{itemize}
    \item \textbf{History:} the mechanism and timing of the injury, whether consciousness was lost, any vomiting or seizures, and current medications, especially anticoagulants and antiplatelet agents
    \item \textbf{Neurological examination:} assessment of conscious level using the Glasgow Coma Scale, pupillary response, limb strength, sensation, coordination, and speech
    \item \textbf{CT scan of the head:} the key investigation to detect bleeding, skull fracture, and brain swelling; it is used routinely in moderate and severe injuries and in mild injuries with risk factors
    \item \textbf{MRI:} may be performed later to characterise diffuse axonal injury or subtle contusions not visible on CT
    \item \textbf{Imaging of the cervical spine:} X-ray or CT to exclude a neck injury, which is assumed present until proven otherwise
    \item \textbf{In-hospital monitoring:} serial neurological observations and, in severe cases, intracranial pressure monitoring to guide management
\end{itemize}

\section*{Management}
\begin{itemize}
    \item \textbf{Mild injury:} initial rest, simple analgesia, and a graduated return to normal activities, together with clear advice about symptoms that should prompt urgent review
    \item \textbf{Observation:} a period of observation in the emergency department or hospital for those at higher risk of deterioration
    \item \textbf{Emergency stabilisation:} airway protection, oxygenation, intravenous fluids, and careful control of blood pressure to prevent secondary injury
    \item \textbf{Surgery:} evacuation of epidural and subdural haematomas, elevation of depressed skull fractures, and decompressive craniectomy for dangerously raised intracranial pressure
    \item \textbf{Intensive care:} management of severe injuries, including control of cerebral oedema, treatment of seizures, and prevention of complications such as pneumonia and venous thromboembolism
    \item \textbf{Rehabilitation:} physiotherapy, occupational therapy, speech and language therapy, and neuropsychological support, ideally through a dedicated brain injury rehabilitation team
    \item \textbf{Follow-up:} monitoring for complications such as post-traumatic epilepsy, and practical support with returning to work, study, sport, and driving
\end{itemize}

\section*{Complications}
\begin{itemize}
    \item Post-concussion syndrome: persistent headache, dizziness, and concentration and memory problems after a mild injury
    \item Post-traumatic epilepsy: seizures that may begin weeks, months, or years after the injury
    \item Permanent neurological deficits such as hemiparesis, aphasia, or visual loss
    \item Cognitive and behavioural change, including impaired memory, impulsivity, and personality alteration
    \item Hydrocephalus from bleeding or infection obstructing cerebrospinal fluid flow
    \item Meningitis or brain abscess, particularly after penetrating injuries or skull-base fractures
    \item Cumulative damage from repeated injuries, which in some athletes leads to chronic traumatic encephalopathy
\end{itemize}

\section*{Prognosis and Outcomes}
Most mild brain injuries resolve fully, although a minority of people experience symptoms that take weeks or months to settle. Recovery from moderate and severe injuries is highly variable and continues over many months or even years: some individuals make substantial functional gains with rehabilitation, while others are left with lasting physical, cognitive, or behavioural disability. Younger age, milder initial injury, and access to prompt specialist care are all associated with better outcomes. Because the injured brain is more vulnerable in the weeks after a concussion, avoiding further head injury is essential, and graduated return-to-sport and return-to-work programmes reduce the risk of re-injury and prolonged symptoms.

\section*{Prevention and Self-Care}
\begin{itemize}
    \item Wear a correctly fitted helmet for cycling, motorcycling, horse riding, skiing, and high-risk or contact sport
    \item Always wear a seatbelt, and never drive under the influence of alcohol or drugs
    \item Reduce fall risk in older adults by removing home hazards, improving lighting, and reviewing medicines that cause dizziness or low blood pressure
    \item Use safety gates, window guards, and age-appropriate restraints around young children
    \item Follow evidence-based return-to-play and return-to-learn guidelines after any concussion
    \item Never move a person with a suspected serious head or neck injury unless it is necessary for safety
    \item Seek medical review after any head injury if symptoms do not settle as expected
\end{itemize}

\section*{Sources and Further Reading}
The following organisations provide reliable, evidence-based information for patients, students, and clinicians.

\begin{itemize}
    \item NHS. \textit{Head injury and concussion: assessment and recovery guidance.} https://www.nhs.uk/conditions/
    \item Centers for Disease Control and Prevention. \textit{Traumatic brain injury and concussion: public health information.} https://www.cdc.gov/
    \item Mayo Clinic. \textit{Traumatic brain injury: overview and patient education.} https://www.mayoclinic.org/diseases-conditions/
    \item MSD Manual Professional Version. \textit{Traumatic brain injury: classification and management.} https://www.msdmanuals.com/
    \item MedlinePlus. \textit{Brain injury and concussion: consumer health information.} https://medlineplus.gov/
    \item World Health Organization. \textit{Global guidance on injury prevention and traumatic brain injury.} https://www.who.int/
\end{itemize}
